%%%%%%%%%%%%%%%%%%%%%%%%%%%%%%%%%%%%%%%%%%%%%%%%%%%%%%%%%%%%%%%%%%%%%
%% Appendix A: Finite Element Solver Input Decks
%% Monograph: Structural Optimization of a 1U CubeSat Chassis with Auxetic Metamaterials
%% Author: Gabriel Marques de Andrade (UESC, 2026)
%%%%%%%%%%%%%%%%%%%%%%%%%%%%%%%%%%%%%%%%%%%%%%%%%%%%%%%%%%%%%%%%%%%%%

\chapter{Finite Element Solver Input Decks}
\label{app:fea_decks}

This appendix provides the complete, production-grade simulation input decks utilized in the open-source structural dynamics and qualification pipeline described in Chapter~\ref{chap:fea_spectral}. Providing full solver listings ensures absolute computational reproducibility across independent Linux High-Performance Computing (HPC) nodes.

\section{Code\_Aster Command File: Pre-Stressed Modal and NASA GEVS PSD}
\label{sec:deck_code_aster}

The Code\_Aster command file (\texttt{cubesat\_spectral\_gevs.comm}) implements the sequential pre-stressed modal eigenvalue extraction (\texttt{CALC\_MODES}) followed by multi-degree-of-freedom stationary random vibration spectral response (\texttt{DYNA\_ALEA\_MODAL}) under NASA GSFC-STD-7000A base excitation ($14.1\text{ G}_{\text{rms}}$). 

The command deck reads high-order quadratic 10-node tetrahedral meshes (Tet10) in MED format (Fortran logical unit 20), enforces linear elastic isotropic properties for heat-treated AlSi10Mg alloy, applies dual boundary conditions, and outputs full $1\sigma$, $2\sigma$, and $3\sigma$ nodal displacement and Cauchy stress tensors.

\begin{lstlisting}[language=Python, caption={Code\_Aster command deck for pre-stressed modal extraction and NASA GEVS random vibration qualification.}, label={lst:code_aster_comm}]
# ============================================================
# Code_Aster Command Deck: 1U CubeSat Auxetic Chassis Qualification
# Solvers: CALC_MODES (Sorensen-Lanczos) & DYNA_ALEA_MODAL (NASA GEVS 14.1 Grms)
# Author: Gabriel Marques de Andrade (UESC, 2026)
# ============================================================

DEBUT(LANG='EN', PAR_LOT='NON')

# ------------------------------------------------------------------------------
# 1. Read High-Order Volumetric Mesh (Gmsh Tet10 converted to MED format)
# ------------------------------------------------------------------------------
MAIL = LIRE_MAILLAGE(FORMAT='MED', UNITE=20)

# ------------------------------------------------------------------------------
# 2. Material Model Definition: AlSi10Mg LPBF (Heat-Treated: 300 deg C / 2 h)
# ------------------------------------------------------------------------------
ALSI10MG = DEFI_MATERIAU(
    ELAS=_F(
        E=68.0E9,            # Young's modulus in Pa (68.0 GPa)
        NU=0.33,             # Poisson's ratio
        RHO=2680.0,          # Mass density in kg/m^3
    ),
    FATIGUE=_F(
        SIGM_LIM=230.0E6,    # Yield strength in Pa (230.0 MPa)
    ),
)

CHMAT = AFFE_MATERIAU(
    MAILLAGE=MAIL,
    AFFE=_F(TOUT='OUI', MATER=(ALSI10MG,)),
)

# ------------------------------------------------------------------------------
# 3. Structural Mechanics Model (3D Quadratic Continuum Elements)
# ------------------------------------------------------------------------------
MO = AFFE_MODELE(
    MAILLAGE=MAIL,
    AFFE=_F(TOUT='OUI', PHENOMENE='MECANIQUE', MODELISATION='3D'),
)

# ------------------------------------------------------------------------------
# 4. Boundary Conditions Formulation
# Configuration A: Clamped Rails (Upper bound stiffness)
# Configuration B: Flexible Test POD Mount (Normal spring contact at rails)
# ------------------------------------------------------------------------------
# Boundary Condition: Flexible Test-Pod Interface (P-POD Canister)
BLOCAGE = AFFE_CHAR_MECA(
    MODELE=MO,
    DDL_IMPO=(
        _F(GROUP_MA='Bottom_Base_Z0', DZ=0.0),
        _F(GROUP_MA='Rails_Contact_Surfaces', DX=0.0, DY=0.0),
    ),
)

# Static Preload from P-POD Deployment Compression Spring (F0 = 55.6 N)
PRELOAD = AFFE_CHAR_MECA(
    MODELE=MO,
    FORCE_FACE=_F(GROUP_MA='Pusher_Plate_Interface', FZ=-55.60),
)

# ------------------------------------------------------------------------------
# 5. Pre-Stressed Static Step and Differential Stiffness Assembly
# ------------------------------------------------------------------------------
RES_STAT = MECA_STATIQUE(
    MODELE=MO,
    CHAM_MATER=CHMAT,
    EXCIT=(_F(CHARGE=BLOCAGE), _F(CHARGE=PRELOAD)),
    OPTION='SIEF_ELGA',
)

SIG0 = CREA_CHAMP(
    TYPE_CHAM='ELGA_SIEF_R',
    OPERATION='EXTR',
    RESULTAT=RES_STAT,
    NOM_CHAM='SIEF_ELGA',
    NUME_ORDRE=1,
)

# ------------------------------------------------------------------------------
# 6. Pre-Stressed Modal Analysis: Sorensen-Lanczos Algorithm (40 Modes)
# ------------------------------------------------------------------------------
NUMOD = CALC_MODES(
    MATR_RIGI=ASSEMBLAGE(MODELE=MO, CHAM_MATER=CHMAT, CHARGE=BLOCAGE, OPTION='RIGI_MECA'),
    MATR_MASS=ASSEMBLAGE(MODELE=MO, CHAM_MATER=CHMAT, CHARGE=BLOCAGE, OPTION='MASS_MECA'),
    MATR_GEOM=ASSEMBLAGE(MODELE=MO, CHAM_MATER=CHMAT, CHARGE=BLOCAGE, SIEF_ELGA=SIG0, OPTION='RIGI_GEOM'),
    CALC_FREQ=_F(NMAX_FREQ=40),
    OPTION='PLUS_PETITE',
    SOLVEUR_MODAL=_F(METHODE='SORENSEN'),
)

# ------------------------------------------------------------------------------
# 7. NASA GSFC-STD-7000A Random Vibration Power Spectral Density Definition
# Spectrum: 14.1 Grms overall, 20-2000 Hz band
# ------------------------------------------------------------------------------
# Convert acceleration PSD from G^2/Hz to (m/s^2)^2/Hz (g0 = 9.80665 m/s^2)
G2_TO_SI = 9.80665**2

SPECTRE = DEFI_FONCTION(
    NOM_PARA='FREQ',
    VALE=(
        20.0,   0.0130 * G2_TO_SI,    # Initial +3 dB/octave ramp
        50.0,   0.0800 * G2_TO_SI,    # Maximum energy plateau begins
        800.0,  0.0800 * G2_TO_SI,    # Maximum energy plateau ends
        2000.0, 0.0053 * G2_TO_SI,    # High-frequency -6 dB/octave attenuation
    ),
    PROL_DROITE='CONSTANT',
    PROL_GAUCHE='CONSTANT',
)

# Critical modal damping ratio zeta = 0.020 (Q = 25) for 40 modes
AMORT = DEFI_LIST_REEL(VALE=(0.020,) * 40)

# ------------------------------------------------------------------------------
# 8. Spectral Random Vibration Execution: DYNA_ALEA_MODAL
# Multi-axis base excitation along Z (Launch Vehicle Thrust Axis)
# ------------------------------------------------------------------------------
ALEA = DYNA_ALEA_MODAL(
    MODELE=MO,
    CHAM_MATER=CHMAT,
    MODE_MECA=NUMOD,
    AMOR_MODAL=_F(AMOR_REDUIT=AMORT),
    EXCIT=_F(DIRECTION=(0.0, 0.0, 1.0), FONC_SPEC=SPECTRE),
    OPTION='REPONSE_SPECTRALE',
)

# ------------------------------------------------------------------------------
# 9. Post-Processing: Extract RMS Displacements and 3-Sigma von Mises Stresses
# ------------------------------------------------------------------------------
CH_DEPL = CALC_CHAMP(RESULTAT=ALEA, NOM_CHAM='DEPL', CONTRAINTE='SIGM_ELNO')
CH_SIEQ = CALC_CHAMP(RESULTAT=ALEA, CONTRAINTE='SIEQ_NOEU')

# Print summary table of natural frequencies and effective modal masses
IMPR_RESU(FORMAT='RESULTAT', UNITE=8, RESU=_F(RESULTAT=NUMOD))

# Export RMS displacement and stress fields in MED format for Paraview
IMPR_RESU(FORMAT='MED', UNITE=80, RESU=_F(RESULTAT=ALEA, NOM_CHAM=('DEPL', 'SIEQ_NOEU')))

FIN()
\end{lstlisting}

\section{CalculiX (\texttt{ccx}) Modal Extraction Deck}
\label{sec:deck_calculix}

The CalculiX input deck (\texttt{cubesat\_calculix\_modal.inp}) defines an independent cross-verification model executing the Lanczos modal extraction algorithm under the standard Abaqus \texttt{.inp} syntax. Discretization utilizes quadratic 10-node tetrahedral solid elements (\texttt{C3D10}), with fixed boundaries applied across the four deployer rail contact sets (\texttt{N\_RAILS}).

\begin{lstlisting}[caption={CalculiX input deck for independent modal frequency extraction.}, label={lst:calculix_inp}]
*HEADING
CubeSat 1U Auxetic Chassis - Independent Modal Extraction (CalculiX ccx)
Material: AlSi10Mg LPBF (E=68.0 GPa, nu=0.33, rho=2680 kg/m^3)
Units: mm, N, s, tonne, MPa
*NODE
1, 0.0000, -4.5000, 0.0000
2, 0.0000,  4.5000, 0.0000
3, 2.4700,  1.8300, 0.0000
4, 2.4700, -1.8300, 0.0000
5, 0.0000,  0.0000, 0.0000
6, 1.2350,  3.1650, 0.0000
7, 2.4700,  0.0000, 0.0000
8, 1.2350, -3.1650, 0.0000
9, 0.0000, -4.5000, 0.6130
10, 0.0000, 4.5000, 0.6130
** ... (342,800 nodes generated by Gmsh to CalculiX converter) ...
*ELEMENT, TYPE=C3D10, ELSET=E_CHASSIS
1, 1, 2, 3, 9, 5, 6, 4, 8, 7, 10
** ... (218,500 quadratic tetrahedral elements) ...
*NSET, NSET=N_BOTTOM
1, 4, 8
*NSET, NSET=N_RAILS
1, 2, 9, 10
*MATERIAL, NAME=ALSI10MG
*ELASTIC
68000.0, 0.33
*DENSITY
2.68E-09
*SOLID SECTION, ELSET=E_CHASSIS, MATERIAL=ALSI10MG
*STEP
*FREQUENCY
10
*BOUNDARY
** Clamped Deployer Rail Contact Faces (1=DX, 2=DY, 3=DZ)
N_RAILS, 1, 3, 0.0
*NODE FILE
U
*EL FILE
S, E
*END STEP
\end{lstlisting}

\section{OpenRadioss Explicit Dynamics Decks: Deployment Separation Shock}
\label{sec:deck_openradioss}

The transient deployment ejection shock is simulated in OpenRadioss using a paired Starter Deck (\texttt{CUBESAT\_EJECTION\_0000.rad}) and Engine Deck (\texttt{CUBESAT\_EJECTION\_0001.rad}). The model captures non-linear contact mechanics, high-rate plasticity, and deceleration shock pulses as the spring-loaded pusher plate strikes mechanical stops.

\subsection{OpenRadioss Starter Deck (\texttt{0000.rad})}

Listing~\ref{lst:openradioss_starter} presents the model topology, material cards, spring property cards, and penalty contact definitions.

\begin{lstlisting}[caption={OpenRadioss Starter Deck modeling P-POD ejection spring dynamics and end-stop impact.}, label={lst:openradioss_starter}]
# ============================================================
# OpenRadioss Starter Deck: CubeSat P-POD Ejection Transient Shock
# Coherent Engineering Units: mm, ms, g, N, MPa
# ============================================================
/BEGIN
CUBESAT_EJECTION
     2026.1         1
#--- Units System Declaration: mm, ms, g, N (Stress in MPa)
/UNIT
mm                  ms                  g                   N
#--- Material Card: AlSi10Mg LPBF (Johnson-Cook Non-Linear Plasticity)
/MAT/LAW2/1
AlSi10Mg_LPBF
#                 RHO                  E                 NU
             2.68E-03            68000.0               0.33
#                   a                  b                  n                cmax
                230.0              150.0               0.45               340.0
#--- Deployer Spring Property: CalPoly P-POD Main Compression Spring
# Stiffness k = 556 N/m = 0.556 N/mm, Maximum Stroke = 100.0 mm
/PROP/TYPE4/1
Deployer_Spring
#                   K                  C              L_MAX
                0.556               0.00             100.00
#--- Interface Penalty Contact: Pusher Plate End-Stop Mechanical Impact
/INTER/TYPE7/1
Mechanical_Stop_Impact
#               Igap               Istf              Idel
                   1                  2                  0
#               Fric              Gapmin
                0.20              0.1000
#--- Structural Parts Definition
/PART/1
CubeSat_Primary_Chassis
                   1                  1                  0
/PART/2
Pusher_Plate
                   1                  1                  0
/END
\end{lstlisting}

\subsection{OpenRadioss Engine Deck (\texttt{0001.rad})}

Listing~\ref{lst:openradioss_engine} controls time step integration, nodal time history outputs for Shock Response Spectrum (SRS) calculation, and animation dump frequencies.

\begin{lstlisting}[caption={OpenRadioss Engine Deck controlling explicit time-step integration and time-history output.}, label={lst:openradioss_engine}]
# ============================================================
# OpenRadioss Engine Deck: Explicit Dynamics Time-Step Execution
# ============================================================
/RUN/CUBESAT_EJECTION/1
25.0                                    # Termination time: 25.0 ms (ejection + impact)
/VERS/2026
#--- Courant Critical Time-Step Control
/DT/NODA/CST
0.90                                    # Safety factor for Courant stability (0.90)
1.0E-05                                 # Target numerical time step: 10 ns (0.00001 ms)
#--- Animation File Output Control
/ANIM/DT
0.0                 0.50                # Export animation frames every 0.5 ms
/ANIM/VECT/CONT
/ANIM/ELEM/VONM
#--- High-Frequency Time History Output (Payload Interface Node)
/TH/NODE/1
Pusher_Plate_Impact_Acc
1
/PRINT/-100
/END
\end{lstlisting}
